\begin{textsample}{NEG Dim 1 – 2010\_c5 – Score 2.00 – 2010\_c5/enade\_2010\_agro\_c3ab28c6eb.txt}  \label{ex:f1_neg_303}
CAPÍTULO 5 DISTRIBUIÇÃO DOS CONCEITOS 5.1 PANORAMA NACIONAL DA DISTRIBUIÇÃO DOS CONCEITOS A Tabela 5.1 apresenta a quantidade e a distribuição de cursos de Agronomia participantes do ENADE/2010, por faixa de conceito e Grande Região.
Observando -se os dados da Tabela 5.1, nota -se que, dos 208 cursos participantes, 52 ( 25,0 \% ) classificaram -se com conceito 3, o valor modal depois da categoria sem conceito ( SC ).
Este não representa, no entanto, o conceito modal em todas as regiões brasileiras.
A região Centro-Oeste, por exemplo, recebeu conceito modal 4, e a Norte, conceito 2, com percentuais de, respectivamente, 23,1 \%, e 23,8 \%.
Dos 208, 60 cursos ( 28,8 \% ) ficaram sem conceito ( SC ).
A região Norte participou com 21 cursos, o que representa 10,1 \% do total nacional.
Nessa região, o conceito 4 foi atribuído a dois cursos dos 21 participantes, equivalente a 9,5 \% do total regional.
Os conceitos 2 ( modal ) e 3 foram atribuídos a cinco ( 23,8 \% ) e a quatro cursos ( 19,0 \% ), respectivamente.
Oito dos cursos ( 38,1 \%, ) não receberam conceito.
Apenas um curso recebeu o conceito 1, com 4,8 \% do total regional.
A região Nordeste participou com 31 cursos, correspondentes a 14,9 \% do total nacional.
Destes, 11 cursos, 35,5 \% em termos regionais, obtiveram conceito 3.
Os conceitos 2 e 4 foram atribuídos a sete cursos para \textbf{cada} um dos conceitos, com percentuais idênticos : 22,6 \%.
O conceito 5 foi atribuído a um curso ( 3,2 \% ).
Nesta região, 16,1 \% dos cursos ( cinco deles ) ficaram sem conceito.
Dos 63 cursos participantes da região Sudeste, 16 ( 25,4 \% ) obtiveram conceito 3, o modal.
O conceito 2 foi atribuído a seis cursos ( 9,5 \% ) e o conceito 4, a 14 ( 22,2 \% ).
Dois cursos ( 3,2 \% ) receberam o conceito 1.
Outros cinco cursos ( 7,9 \% ) receberam conceito 5.
Os 20 cursos restantes ( 31,8 \% ) ficaram sem conceito.
Quase 1/4 ( 24,1 \%, correspondendo a 13 cursos ) dos 54 participantes na região Sul recebeu conceito 3, o conceito modal.
Nessa região, quatro cursos receberam conceito 2 ( 7,4 \% ) e 11 cursos ( 20,4 \% ) receberam conceito 4.
Nove cursos, equivalentes a 16,7 \%, receberam conceito 5.
Ficaram sem conceito, 29,5 \% do total regional.
Na região Centro-Oeste, o conceito 1 foi obtido por dois cursos ( 5,1 \% ), o conceito 2 por seis deles ( 15,4 \% ), o 3 foi atribuído a oito cursos ( 20,5 \% ) e o conceito 5, a três cursos ( 7,7 \% ).
Esta região teve o conceito 4 como o modal, como já mencionado, destacando -se nove cursos, 23,1 \% do total para essa região.
Do total de cursos da região Centro-Oeste, 11 ( 28,2 \% ) ficaram sem conceito. 5.2 CONCEITOS POR CATEGORIA ADMINISTRATIVA E POR GRANDE REGIÃO A Tabela 5.2 apresenta a distribuição dos cursos participantes do ENADE/2010 de Agronomia, por Categoria Administrativa, de acordo com os conceitos por eles alcançados, segundo as Grandes Regiões brasileiras.
Dos 208 cursos participantes, 120 ( 57,7 \% ) correspondiam a instituições públicas e 88 ( 42,3 \% ), a IES privadas.
De acordo com as informações da Tabela 5.2, as instituições públicas foram as que apresentaram, em termos nacionais, a maior proporção de cursos com conceito 5, equivalente a 15,8 \% ( 19 cursos ) do total dessa categoria.
Dos 120 cursos de instituições públicas participantes, o conceito 3 foi atribuído a 19 deles, e o conceito 4, a 34 cursos, sendo este último o valor modal.
Entre os demais participantes, onze cursos receberam conceito 2, e dois deles o conceito 1.
Ficaram sem conceito 35 cursos.
Na rede privada, o conceito modal foi 3, com 33 cursos.
Entre os demais cursos participantes, 17 receberam conceito 2 e nove, o conceito 4.
O conceito 1 foi atribuído a quatro cursos.
Nesta categoria, 25 cursos ficaram sem conceito.
Na análise por região, observa -se que na região Norte as instituições públicas participaram com 14 cursos, dos quais cinco ficaram sem conceito.
O conceito modal para as instituições públicas na região foi distribuído entre 2 e 3, com três cursos \textbf{cada}.
Os demais foram avaliados com conceito 4 ( dois cursos ) e conceito 5 ( um curso ).
As instituições privadas participaram com sete cursos, dos quais três ficaram sem conceito.
O conceito modal para esta Categoria Administrativa, na região Norte, foi 2, com dois cursos.
Os conceitos 1 e 3 foram atribuídos a um curso \textbf{cada}.
Na região Nordeste, a rede pública concentrou 29 dos 31 cursos participantes, 93,5 \% do total da região.
Entre os que receberam conceito nesta Categoria Administrativa na região Nordeste, predominaram os cursos com conceito 2 ( 11 cursos ), concentrados em IES públicas.
Neste tipo de instituição, os conceitos 1 e 3 ficaram, \textbf{cada} qual, com seis cursos.
Os demais cursos foram avaliados com conceito 4 ( um curso ) e conceito 5 ( cinco cursos ).
As instituições privadas, nesta região, participaram com apenas dois cursos ( 6,5 \% ).
O conceito modal para esta Categoria Administrativa ficou entre 2 e 4, com um curso \textbf{cada}.
Na região Sudeste, a proporção de cursos da rede privada, 58,7 \%, foi mais elevada do que nas IES públicas, correspondendo a 37 dos 63 cursos participantes.
Em IES privadas, na região Sudeste, o conceito modal foi 3 ( 15 cursos ).
Os demais cursos foram avaliados com conceito 1 ( dois cursos ), conceito 2 ( seis cursos ) e conceito 4 ( quatro cursos ).
Nesta combinação de Categoria Administrativa e Grande Região, dez cursos não receberam conceito.
Entre os 26 cursos em instituições públicas na região, a categoria modal foi a que recebeu o conceito 4 ( dez cursos ).
Os demais foram avaliados com conceito 3 ( um curso ) e conceito 5 ( cinco cursos ).
Nesta categoria na região, dez cursos ficaram sem conceito.
Tanto as instituições públicas quanto as privadas dividiram em igual medida os 57 cursos participantes da região Sul.
Em \textbf{cada} uma dessas Categorias Administrativas, houve 27 cursos participantes, 50,0 \% do total regional para \textbf{cada} tipo de IES.
Nas instituições públicas, o conceito modal ficou dividido entre os conceitos 4 e 5 ( com nove cursos \textbf{cada} ).
Um curso obteve conceito 1 e outro, o conceito 3.
Na combinação de Categoria Administrativa e Grande Região, sete cursos não receberam conceito.
Nas instituições privadas, nesta região, o conceito modal foi 3, com 12 cursos.
Excluindo -se os nove cursos que ficaram sem conceito, os restantes se distribuíram pelo conceito 2 ( quatro cursos ) e conceito 4 ( dois cursos ).
Na região Centro-Oeste, 24 dos 39 cursos participantes eram de instituições públicas ( 61,5 \% em termos regionais ).
Destes, sete concentraram -se no conceito modal, 4, e oito ficaram sem conceito.
Os demais receberam conceito 1 ( um curso ), conceito 2 ( dois cursos ), 3 ( três cursos ) e 5 ( três cursos ).
Na rede privada, dos 15 cursos na região Centro-Oeste, cinco foram avaliados com conceito modal, 3.
Ficaram sem conceito três cursos.
Os outros cursos foram avaliados com os conceitos 1 ( um curso ), 2 ( quatro cursos ) e conceito 4 ( dois cursos ). 5.3 CONCEITOS POR ORGANIZAÇÃO ACADÊMICA E POR GRANDE REGIÃO Na Tabela 5.3 encontra -se a distribuição dos conceitos atribuídos aos cursos participantes do ENADE/2010 na Área de Agronomia, por Organização Acadêmica, segundo as Grandes Regiões brasileiras.
Dos 208 cursos participantes, 143 eram oferecidos em Universidades, 14 em Centros Universitários e, os demais 51, em Faculdades.
De acordo com os dados apresentados, todos os cursos avaliados com conceito 5 eram vinculados a Universidades ( o total dos 19 cursos com este conceito ).
Esse tipo de Organização Acadêmica teve o conceito 4 como modal, com 39 cursos.
Além dos outros 39 cursos que ficaram sem conceito, os demais foram avaliados com conceito 1 ( um curso ), conceito 2 ( 15 cursos ) e conceito 3 ( 30 cursos ).
Entre os cursos em Centros Universitários, o conceito modal foi 3, com cinco cursos, sendo que o mesmo número de cursos ficou sem conceito.
Os outros cursos neste tipo de Organização Acadêmica receberam conceito 1 ( um curso ), 2 ( dois deles ) e conceito 4 ( apenas um curso ).
Nas Faculdades, dos 51 cursos, 16 ficaram sem conceito e 17 receberam o mesmo conceito modal dos Centros Universitários, 3.
Além dos cursos que receberam o conceito 1 ( quatro cursos ), os demais receberam conceito 2 ( 11 cursos ) e conceito 4 ( três cursos ).
Considerando -se separadamente as regiões brasileiras, verifica -se que, na região Norte, as Universidades concentraram 14 dos 21 cursos participantes.
Dos cursos em Universidades, cinco ficaram sem conceito e os restantes foram avaliados com conceito 2 ( três cursos ) e conceito 3 ( três cursos ).
Essa quantidade igual de cursos indica terem sido 2 e 3 os conceitos modais.
Receberam o conceito 4 dois cursos, e o conceito 5, foi a avaliação de apenas um deles.
Os Centros Universitários da região Norte foram representados por dois cursos, sendo um com o conceito modal 2 e um sem conceito.
As Faculdades participaram com cinco cursos.
Destes, dois ficaram sem conceito, sendo que os três cursos restantes obtiveram, um em \textbf{cada}, os conceitos 1, 2 e 3.
Na região Nordeste, as Universidades participaram com 29 cursos dos 31 na Área de Agronomia.
Entre os cursos em Universidades, cinco ficaram sem conceito.
O conceito modal foi 3, com 11 cursos.
Os demais receberam conceito 2 ( seis cursos ), conceito 4 ( também seis ) e conceito 5 ( um curso ).
Na região Nordeste não houve participação de cursos em Centros Universitários.
As Faculdades, por sua vez, foram representadas por apenas dois cursos, sendo que um recebeu conceito 2 e o outro, conceito 4.
Nesse caso, o modal ficou entre 2 e 4.
Na região Sudeste, as Universidades concentraram 33 dos 63 cursos da região.
Entre os cursos em Universidades na região, o conceito modal foi 4, com 12 cursos, e 11 ficaram sem conceito.
Os demais receberam conceito 3 ( cinco cursos ) e conceito 5 ( também cinco cursos ).
Os Centros Universitários participaram com sete cursos na região Sudeste, dos quais três obtiveram o conceito modal, 3, e três ficaram sem conceito.
O restante recebeu conceito 1.
As Faculdades foram representadas por 23 cursos na região Sudeste, que se distribuíram nos conceitos 1 ( um curso ), 2 ( seis cursos ), 3 ( oito cursos, o conceito modal ), e 4 ( dois cursos ).
Os cursos sem conceito foram seis.
Dos 54 cursos da região Sul, 42 eram de Universidades, para os quais o conceito modal foi 4 ( com 11 cursos ).
Nesse tipo de Organização, dez cursos ficaram sem conceito e os demais receberam os conceitos 1 ( apenas um curso ), 2 ( quatro cursos ) e conceito 3 ( sete cursos ).
Os Centros Universitários da região Sul tiveram apenas um dos dois cursos participantes no conceito modal, o conceito 3.
O outro curso não recebeu conceito.
Foram dez os cursos vinculados a Faculdades na região Sul, e destes, cinco ficaram sem conceito.
Os cinco restantes receberam conceito 3.
Na região Centro-Oeste, 25 dos 39 cursos eram de Universidades.
Destes, oito cursos ficaram sem conceito e o conceito modal foi 4, com oito cursos.
Os cursos além Os cursos além destes foram avaliados com conceito 2 ( dois cursos ), 3 ( quatro deles ) e conceito 5 ( três cursos ).
Os Centros Universitários da região Centro-Oeste contaram com três cursos, cuja avaliação resultou em uma distribuição pelos conceitos 2, 3 e 4.
Dos 11 cursos em Faculdades na região Centro-Oeste, três ficaram sem conceito.
Dois deles receberam conceito 1, sendo que o conceito modal ficou distribuído entre 2 e 3, com três cursos \textbf{cada}.

% matched lemmas: cada
\end{textsample}
